\section{Deflated logarithms, contour floors and signed integration}
\label{sec:deflation}

The generic logarithmic envelope uses no real-boundary modulus floor.
When a certified local divisor is available, one can instead integrate
its singular logarithms explicitly and approximate the remaining regular
factor. The following construction is inherited from the assessed
deflation note and its later conditional contour design
\cite{OGDeflation,OGDeflationAssessment,OGContourDesign}.

\subsection{Stable factor acquisition by divided differences}
On a real patch \(I_v\), suppose the relevant actual roots of \(H_0\)
are the multiset \(z_1,\ldots,z_n\), and define the removable quotient
\[
 U_v(z)=H_0(z)\Big/\prod_{j=1}^n(z-z_j).
\]
Repeated roots are repeated in the list. Deflate \(H_0\), whose entire
domain is available, before multiplying by \(D\); the meromorphic
extension of \(F\) has a pole at \(-\ii\).

For a holomorphic \(f\), its divided difference has the simplex formula
\begin{equation}\label{eq:divided-simplex}
 f[z_1,\ldots,z_n,x]
 =\int_{\Delta_n}f^{(n)}\left(t_xx+\sum_{j=1}^nt_jz_j\right)\dd t,
 \qquad \operatorname{vol}(\Delta_n)=1/n!.
\end{equation}
It follows by iterating the fundamental theorem of calculus; continuity
in the knots supplies repeated-node values. When the knots are the
actual roots, with their multiplicities, the interpolation polynomial
vanishes and \eqref{eq:divided-simplex} equals \(U_v(x)\).

Suppose certified matching gives approximate knots \(w_j\) with
\(|z_j-w_j|\le\rho_j\), and let \(\widehat H\) approximate \(H_0\)
on the full convex domain of the knots and patch. Define
\(\widehat U_v(x)=\widehat H[w_1,\ldots,w_n,x]\).
If \(\epsilon_k\) bounds the error in derivative \(k\) of \(H_0\), and
\(\widehat M_k\) bounds \(|\widehat H^{(k)}|\), then
\begin{equation}\label{eq:deflated-errors}
 \|U_v^{(j)}-\widehat U_v^{(j)}\|_{I_v}
 \le e_j:=\frac{j!\epsilon_{n+j}}{(n+j)!}
       +\frac{j!\widehat M_{n+j+1}}{(n+j+1)!}\sum_{l=1}^n\rho_l,
 \quad j=0,1,2.
\end{equation}
Differentiate the simplex integral: the factor is \(t_x^j\).
Moving a knot adds \(t_l\rho_l\) and one derivative. The simplex
moments are \(j!/(n+j)!\) and \(j!/(n+j+1)!\), respectively.
This proves the estimate without division by a root separation.
Approximate knots require derivatives through order \(n+3\).
The function \(\widehat U_v\) is a divided difference, not
\(\widehat H/\prod(z-w_j)\); the latter would leave spurious poles.

If \(|\widehat U_v|\ge\mu_v>e_0\) on the real patch, integration of
the derivative of the real logarithm along a segment of moduli gives
\begin{equation}\label{eq:factor-log-error}
 |\log|U_v|-\log|\widehat U_v||\le b_v:=\frac{e_0}{\mu_v-e_0}.
\end{equation}
With \(W_v=\int_{I_v}K(x-X)\dd x\), the factor error is
\(E_{\rm fac}=\sum_vW_vb_v\), and \(\sum_vW_v\le1/2\).
The floor is an additional premise, not an automatic consequence of
divided differences.

\subsection{Paying root uncertainty through real zeros}
For an interval of length \(L\), two complex points within distance
\(\rho\) satisfy the root-log estimate from the deflation source
\cite[Section 3]{OGDeflation}
\begin{equation}\label{eq:root-log-error}
 \int_I|\log|x-z|-\log|x-w||\dd x
 \le A(L,\rho):=4\rho\left[1+\log\left(1+\frac{L}{2\rho}\right)\right],
 \quad A(L,0)=0.
\end{equation}
For a direct proof, the triangle inequality gives
\[
 |\log|x-z|-\log|x-w||
 \le\log\left(1+\frac{\rho}{\min(|x-z|,|x-w|)}\right)
 \le\sum_{a\in\{\Re z,\Re w\}}\log(1+\rho/|x-a|).
\]
For each center the integral is largest when the interval is centered
there, because the integrand is even and decreases with distance. Its
value is
\[
 2\int_0^{L/2}\log(1+\rho/t)\dd t
 =L\log(1+2\rho/L)+2\rho\log(1+L/(2\rho))
 \le2\rho[1+\log(1+L/(2\rho))].
\]
Adding the two centers proves \eqref{eq:root-log-error}, including real
roots by integrable limiting values. No simple-zero slope is used.
Since \(K\le1/16\), one may charge
\[
 E_{\rm roots}=\frac1{16}\sum_{v,j}A(|I_v|,\rho_{vj}).
\]
Repeated uses of a root in different patches are charged again.
This is stability in root location. Source perturbations can move
roots by a square root, so it is not a Lipschitz theorem in source error.

\subsection{Product integration with exact singular primitives}
Set \(q_v(x)=\log|D(x)|+\log|\widehat U_v(x)|\) and assume
\(|q_v|\le H_v\), \(|q_v''|\le B_v\). On panels of maximum width
\(\Delta\), linearly interpolate both \(K\) and \(q_v\), denoting
the interpolants by \(K_h,q_h\). Use
\begin{equation}\label{eq:product-integration}
 Q_h=\sum_v\int_{I_v}K_h(x)
       \left(q_h(x)+\sum_j\log|x-w_{vj}|\right)\dd x.
\end{equation}
The regular part is polynomial integration. For a root with real part
\(a_j\) and \(b_j=|\Im w_j|\), put \(t=x-a_j\). The two singular
primitives needed for the linear kernel are
\begin{align}
 A_0(t,b)&=\tfrac t2\log(t^2+b^2)-t+b\arctan(t/b),\nonumber\\
 A_1(t,b)&=\tfrac14[(t^2+b^2)\log(t^2+b^2)-(t^2+b^2)].
 \label{eq:root-primitives}
\end{align}
Their derivatives are \(\log\sqrt{t^2+b^2}\) and
\(t\log\sqrt{t^2+b^2}\). Continuous limits at \(b=0\) include roots
at panel endpoints; there is no exclusion of multiple real roots.

An explicit common bound for the unweighted absolute integrals of the
regular part and singular logarithms is
\[
 \mathcal S:=\sum_v\left[|I_v|H_v+
       \sum_j\{|I_v|\log^+D_{vj}+2\}\right],\qquad
 D_{vj}=\max_{x\in I_v}|x-w_{vj}|.
\]
The positive logarithm uses its maximum; the integral of the negative
logarithm is at most \(\int_{-1}^1-\log|t|\dd t=2\).
Linear interpolation has uniform error at most \(\Delta^2/8\) times
the second-derivative bound, and \(\|K''\|_\infty\le a^2/16\).
Split the regular integrand as
\(Kq_v-K_hq_h=K(q_v-q_h)+(K-K_h)q_h\).
Linear interpolation preserves the bound \(|q_h|\le H_v\), so the same
\(\mathcal S\) also bounds the integral of the interpolated regular
part plus the singular logarithms. Consequently
\begin{equation}\label{eq:product-error-safe}
 E_{\rm quad}\le\frac{\Delta^2}{8}
       \left(\sum_v W_v B_v+\frac{a^2}{16}\mathcal S\right)
 \le\Delta^2\left(\frac{B_*}{16}+\frac{a^2\mathcal S}{128}\right),
 \qquad B_*=\max_v B_v.
\end{equation}
Only the first term uses exact Poisson weights \(W_v\). Numerical
node errors, when later multiplied by the interpolated kernel, instead
use its explicitly computable weights \(\int K_h\). This distinction
avoids an unproved assumption that interpolation preserves Poisson mass.
For fixed paid constants the panel-width cost is of order
\(E_{\rm quad}^{-1/2}\); the source, count and factor costs remain.

\subsection{Contour acquisition of the factor floor}
Suppose a closed disk of radius \(A\), centered at real \(c\), contains
exactly \(n\) roots with multiplicity, all at distance at most \(s<A\)
from \(c\), and
\(0<m\le|H_0|\le M\) on its boundary. Deflating all these roots gives
a holomorphic, zero-free function \(U\) on the closed disk. The maximum
principle applied to \(U\) and \(1/U\) proves
\begin{equation}\label{eq:contour-factor-floor}
 \mu_0:=\frac{m}{(A+s)^n}\le|U(z)|\le
 V_0:=\frac{M}{(A-s)^n}.
\end{equation}
Thus a paid boundary margin and complete root count can replace a
direct search for a small real-axis denominator. They cannot be omitted.
For a polynomial source model \(p\), Rouch\'e's theorem on the outer
boundary, together with disjoint inner cluster counts summing to the
outer count, supplies the required completeness. A list of located
roots without this count does not rule out missed roots in the annulus.

If \eqref{eq:deflated-errors} holds on the entire disk and
\(e_0<\mu_0\), the divided-difference model \(\widetilde U\) is zero-free
there and lies between \(\mu=\mu_0-e_0\) and \(V=V_0+e_0\).
For a real subinterval \(|x-c|\le\ell<A\), harmonicity of
\(\log|\widetilde U|\) gives
\begin{equation}\label{eq:factor-curvature}
 |(\log|\widetilde U|)''(x)|
 \le\frac{2\log(V/\mu)}{(A-\ell)^2}.
\end{equation}
For example, on a circle of radius \(r<A-\ell\), subtract the midpoint
of its allowed logarithmic range. The second Fourier coefficient of
this bounded harmonic function gives a bound
\(4\log(V/\mu)/(\pi r^2)\le2\log(V/\mu)/r^2\); let \(r\uparrow A-\ell\).
On the real axis, direct differentiation of \eqref{eq:F} gives
\[
 (\log|D|)''=\Re\{-1/(2s^2)+\psi'(s/2)/16\},\qquad
 |(\log|D|)''|\le13/4.
\]
Indeed \(|s|\ge1/2\) and the trigamma series at \(\Re(s/2)=1/4\)
is bounded by \(16+\int_0^\infty(t+1/4)^{-2}\dd t=20\).
Together these bounds supply the \(B_v\) needed in
\eqref{eq:product-error-safe}.

All derivatives used above can be acquired from the original source.
For a degree-\(L\) Taylor polynomial about \(c\), if coefficient
errors are \(\delta_j\), finite-source derivative bound is
\(M_{L+1,A}^{\rm fin}\), and the complete omitted-source derivative
allowance is \(E_{{\rm src},k}\), then
\begin{equation}\label{eq:taylor-source-error}
 \epsilon_k\le E_{{\rm src},k}
 +\sum_{j=k}^L(j)_k\delta_jA^{j-k}
 +M_{L+1,A}^{\rm fin}\frac{A^{L+1-k}}{(L+1-k)!}.
\end{equation}
The moment integral of \(|u|^k\ee^{A|u|}\) against the arithmetic and
radial tails proves the first term. Taylor's integral remainder proves
the last. Increasing arithmetic precision does not remove either of
these analytic remainders.

\subsection{The signed error and an identical-information comparison}
Let a numerical enclosure of \(Q_h-P_Q+G(X,1)\) have center \(S_c\)
and radius \(e_{\rm num}\). Combining the preceding estimates yields
\begin{equation}\label{eq:deflation-ledger}
 |J_0(X)-S_c|\le\mathcal R
 :=\mathcal T+\beta_Q+E_{\rm fac}+E_{\rm roots}+E_{\rm quad}+e_{\rm num}.
\end{equation}
Consequently \(J_{\rm up}=S_c+\mathcal R\) satisfies
\(0\le J_{\rm up}-J_0\le2\mathcal R\). Small acquisition error is
separate from the strict value test \(J_{\rm up}<m(d,r)\).

There is also a useful scoped comparison with direct zero information.
Suppose a complete collection of rectangular cells
\(C_i=[l_i,u_i]\times[b_i,c_i]\) carries integer counts \(n_i\).
For each cell in the upper strip, the independent Green maximum is
\[
 U_i=G(\operatorname{dist}(X,[l_i,u_i]),\max(c_i,0)).
\]
If an occupied cell intersects the target
\([X-r,X+r]+\ii[d,1]\), monotonicity of \(G\) gives
\(U_i\ge G(r,d)\ge m(d,r)\). Therefore a strict pass of
\(\sum_i n_iU_i+T_G<m(d,r)\), with \(T_G\ge0\) a complete tail,
already implies that every occupied cell is disjoint from the target.
The same data thus directly exclude zeros there. This proves dominance
only for independent rectangular maxima with identical complete cell
information. It is neither a cost theorem nor an obstruction to coupled
signed estimates, different information, or reuse of an acquired observable.
